\section{Билет 52. Распределение Максвелла молекул по проекциям скорости (основные идеи, без вывода). График функции распределения.}

\begin{definition}
\textbf{Распределение Максвелла} описывает вероятность того, что молекула идеального газа, находящегося в термодинамическом равновесии при температуре $T$, имеет определённую проекцию скорости на заданную ось (например, $v_x$).
\par\noindent\textbf{Основные идеи, лежащие в основе распределения:}
\begin{enumerate}
    \item Движение молекул хаотично и изотропно: все направления равноправны.
    \item Проекции скорости $v_x$, $v_y$, $v_z$ --- статистически независимые случайные величины.
    \item Распределение должно зависеть только от кинетической энергии молекулы $E = \frac{1}{2} m v^2$ и температуры $T$.
    \item При $T \to 0$ все молекулы должны иметь нулевую скорость; при $T \to \infty$ распределение должно ``размазываться''.
\end{enumerate}
\end{definition}

\begin{definition}
Функция распределения Максвелла по одной проекции скорости $v_x$ имеет вид:
\begin{equation}
    \varphi(v_x) = \left( \frac{m}{2\pi k T} \right)^{1/2} \exp\left( -\frac{m v_x^2}{2 k T} \right), \label{eq:maxwell_vx_52}
\end{equation}
где:
\begin{itemize}
    \item $m$ --- масса одной молекулы;
    \item $k = 1.38 \cdot 10^{-23}$ Дж/К --- постоянная Больцмана;
    \item $T$ --- абсолютная температура.
\end{itemize}
Вероятность того, что проекция скорости молекулы находится в интервале $(v_x, v_x + dv_x)$, равна:
\begin{equation}
    dP = \varphi(v_x) \, dv_x.
\end{equation}
\end{definition}

\begin{property}[Свойства функции $\varphi(v_x)$]
\begin{enumerate}
    \item \textbf{Чётность:} $\varphi(-v_x) = \varphi(v_x)$ --- распределение симметрично относительно нуля, так как направления ``влево'' и ``вправо'' равновероятны.
    \item \textbf{Максимум:} достигается при $v_x = 0$:
    \begin{equation}
        \varphi_{\max} = \varphi(0) = \left( \frac{m}{2\pi k T} \right)^{1/2}.
    \end{equation}
    \item \textbf{Асимптотика:} при $|v_x| \to \infty$ функция экспоненциально убывает: $\varphi(v_x) \to 0$.
    \item \textbf{Нормировка:} $\int_{-\infty}^{+\infty} \varphi(v_x) \, dv_x = 1$.
    \item \textbf{Среднее значение:} $\langle v_x \rangle = 0$ (из симметрии).
    \item \textbf{Средний квадрат:} $\langle v_x^2 \rangle = \frac{kT}{m}$, что согласуется с $\langle E_{\text{пост}} \rangle = \frac{3}{2} kT$.
\end{enumerate}
\end{property}

\begin{remark}
\textbf{График функции распределения $\varphi(v_x)$.}
\begin{itemize}
    \item Кривая имеет форму ``колокола'' (гауссианы), симметричного относительно оси $v_x = 0$.
    \item При увеличении температуры $T$ кривая становится шире и ниже: максимум уменьшается как $1/\sqrt{T}$, а ``разброс'' скоростей растёт.
    \item При увеличении массы молекулы $m$ кривая сужается и становится выше: тяжёлые молекулы имеют меньший разброс скоростей при той же температуре.
    \item Площадь под кривой всегда равна 1 (условие нормировки).
\end{itemize}
\end{remark}

\begin{remark}
\textbf{Физическая интерпретация.} Распределение Максвелла показывает, что:
\begin{enumerate}
    \item Наиболее вероятно, что проекция скорости молекулы близка к нулю.
    \item Вероятность очень больших скоростей экспоненциально мала, но не равна нулю --- это объясняет, например, испарение жидкостей при температурах ниже точки кипения (наиболее быстрые молекулы покидают поверхность).
    \item Распределение не зависит от давления и объёма газа --- только от температуры и массы молекул.
\end{enumerate}
\end{remark}
